% Part: sets-functions-relations
% Chapter: sets
% Section: unions-and-intersections

\documentclass[../../../include/open-logic-section]{subfiles}

\begin{document}

\olfileid{sfr}{set}{uni}
\olsection{Gabungan lan Irisan}

\begin{explain}
Ing \olref[sfr][set][bas]{sec}, kita ngenalake definisi himpunan
kanthi abstraksi, yaiku definisi awujud $\Setabs{x}{\phi(x)}$.
Ing kene kita nggunakake sawijining sipat~$\phi$, lan sipat iki
bisa nyebut himpunan kang wis ditetepake sadurunge. Tuladhane,
yen $A$ lan~$B$ iku himpunan, himpunan
$\Setabs{x}{x \in A \lor x \in B}$ ngemot kabeh objek kang
dadi !!{element}s saka $A$ utawa~$B$, yaiku himpunan kang
nggabungake !!{element}s saka $A$ lan~$B$. Iki bisa digambarake
kaya ing \olref{fig:union}, kanthi dhaerah kang disorot nuduhake
!!{element}s loro himpunan $A$ lan~$B$ bebarengan.

\begin{figure}
  \olasset{assets/diagrams/union.tikz}
  \caption{Gabungan $A \cup B$ saka rong himpunan yaiku himpunan
   !!{element}s saka $A$ bebarengan karo anggota saka~$B$.}
  \ollabel{fig:union}
\end{figure}

Operasi tumrap himpunan iki---nggabungake himpunan---migunani banget
lan kerep digunakake, mula kita menehi jeneng formal lan lambang.
\end{explain}

\begin{defn}[Gabungan]
\emph{Gabungan} rong himpunan $A$ lan $B$, ditulis $A \cup B$,
yaiku himpunan kabeh bab kang dadi !!{element}s saka $A$,
$B$, utawa loro-lorone.
\[
A \cup B = \Setabs{x}{x \in A \lor x \in B}
\]
\end{defn}

\begin{ex}
Amarga cacah pangulangan !!{element}s ora dadi prakara,
gabungan rong himpunan kang duwe !!a{element} kang padha
mung ngemot !!{element} kasebut sapisan, tuladhane
$\{ a, b, c\} \cup \{ a, 0, 1\} = \{a, b, c, 0, 1\}$.

Gabungan himpunan lan salah siji subhimpunane iku persis himpunan
kang luwih gedhe: $\{a,b, c \} \cup \{a \} = \{a, b, c\}$.

Gabungan himpunan karo himpunan kosong padha karo himpunan kasebut:
$\{a,b, c \} \cup \emptyset = \{a, b, c \}$.
\end{ex}

\begin{prob}
Buktekna yen $A \subseteq B$, mula $A \cup B = B$.
\end{prob}

\begin{explain}
Kita uga bisa ngrembug operasi kang "dual" karo gabungan.
Operasi iki mbentuk himpunan kabeh !!{element}s kang dadi
!!{element}s saka~$A$ lan uga !!{element}s saka~$B$.
Operasi iki diarani \emph{irisan}, lan bisa digambarake kaya
ing \olref{fig:intersection}.
\begin{figure}
  \olasset{assets/diagrams/intersection.tikz}
  \caption{Irisan $A \cap B$ saka rong himpunan yaiku himpunan
    !!{element}s kang padha-padha diduweni loro himpunan kasebut.}
  \ollabel{fig:intersection}
\end{figure}
\end{explain}

\begin{defn}[Irisan]
\emph{Irisan} rong himpunan $A$ lan $B$, ditulis $A \cap B$,
yaiku himpunan kabeh bab kang dadi !!{element}s saka $A$ lan~$B$
bebarengan.
\[
A \cap B = \Setabs{x}{x \in A \land x \in B}
\]
Rong himpunan diarani \emph{pisah} yen irisane kosong.
Tegese, loro himpunan iku ora duwe !!{element}s kang padha.
\end{defn}

\begin{ex}
Yen rong himpunan ora duwe !!{element}s kang padha, irisane kosong:
$\{ a, b, c\} \cap \{ 0, 1\} = \emptyset$.

Yen rong himpunan duwe !!{element}s kang padha, irisane yaiku
himpunan kabeh anggota kang padha mau:
$\{a, b, c \} \cap \{a, b, d \} = \{a, b\}$.

Irisan himpunan karo salah siji subhimpunane iku persis
himpunan kang luwih cilik:
$\{a, b, c\} \cap \{a, b\} = \{a, b\}$.

Irisan saben himpunan karo himpunan kosong iku kosong:
$\{a, b, c \}\cap \emptyset = \emptyset$.
\end{ex}

\begin{prob}
Buktekna kanthi tliti miturut paugeran matematis yen $A \subseteq B$,
mula $A \cap B = A$.
\end{prob}

\begin{explain}
Kita uga bisa mbentuk gabungan utawa irisan luwih saka rong
himpunan. Cara kang prasaja lan trep kanggo ngrembug iki kanthi
umum yaiku mangkene: umpamakna kabeh himpunan kang arep digabungake
(utawa dijupuk irisane) diklumpukake dadi siji himpunan.
Gabungan kabeh himpunan wiwitan banjur bisa ditemtokake minangka
himpunan kabeh objek kang kalebu ing paling ora siji
!!{element} himpunan kasebut, lan irisane minangka himpunan kabeh
objek kang kalebu ing saben !!{element} himpunan kasebut.
\end{explain}

\begin{defn}
Yen $A$ iku himpunan kang ngemot himpunan-himpunan,
mula $\bigcup A$ yaiku himpunan !!{element}s saka
!!{element}sne~$A$:
\begin{align*}
\bigcup A & = \Setabs{x}{x \text{ kalebu ing !!a{element} saka } A},
\text{ yaiku,}\\
& = \Setabs{x}{\text{ana } B \in A
  \text{ kang ndadekake } x \in B}
\end{align*}
\end{defn}

\begin{defn}
Yen $A$ iku himpunan kang ngemot himpunan-himpunan,
mula $\bigcap A$ yaiku himpunan objek kang padha-padha diduweni
kabeh anggota~$A$:
\begin{align*}
\bigcap A & = \Setabs{x}{x \text{ kalebu ing saben !!{element} saka } A},
\text{ yaiku,}\\
 & = \Setabs{x}{\text{kanggo kabeh } B \in A, x \in B}
\end{align*}
\end{defn}

\begin{ex}
Umpamakna $A = \{ \{ a, b \}, \{ a, d, e \}, \{ a, d \} \}$.
Mula $\bigcup A = \{ a, b, d, e \}$ lan $\bigcap A = \{ a \}$.
\end{ex}
\begin{prob}
	Tuduhna yen $A$ iku himpunan lan $A \in B$, mula $A \subseteq \bigcup B$.
\end{prob}

Bab kang padha uga bisa ditindakake kanggo urutan himpunan $A_1$, $A_2$, \dots
\begin{align*}
\bigcup_i A_i & = \Setabs{x}{x \text{ kalebu ing salah siji } A_i}\\
\bigcap_i A_i & = \Setabs{x}{x \text{ kalebu ing saben } A_i}.
\end{align*}

Yen ana \emph{indeks} himpunan, yaiku sawijining himpunan $I$
kang ndadekake kita ngrembug $A_i$ kanggo saben $i \in I$,
kita uga bisa nganggo cekakan iki:
\begin{align*}
	\bigcup_{i \in I} A_i & = \bigcup \Setabs{A_i }{i \in I}\\
	\bigcap_{i \in I} A_i & = \bigcap\Setabs{A_i}{i \in I}
\end{align*}

Pungkasane, kita bisa ngrembug himpunan kabeh !!{element}s
ing~$A$ kang ora ana ing~$B$. Iki bisa digambarake kaya
ing \olref{difference}.

\begin{figure}
  \olasset{assets/diagrams/difference.tikz}
  \caption{Selisih $A \setminus B$ saka rong himpunan yaiku
    himpunan !!{element}s saka~$A$ kang ora uga dadi
    !!{element}s saka~$B$.}
  \ollabel{difference}
\end{figure}

\begin{defn}[Selisih]
\emph{Selisih himpunan}~$A \setminus B$ yaiku himpunan kabeh
!!{element}s saka $A$ kang ora uga dadi !!{element}s saka~$B$,
yaiku,
\[
A\setminus B = \Setabs{x}{x\in A \text{ lan } x \notin B}.
\]
\end{defn}

\begin{prob}
	Buktekna yen $A \subsetneq B$, mula $B \setminus A \neq \emptyset$.
\end{prob}

\end{document}
